\chapter*{Abstract}
\addcontentsline{toc}{chapter}{Abstract}

This monograph presents Metatime/TIR as an exploratory, low-parameter
phenomenological ansatz relating geometric-phase and information-theoretic
structures to particle, flavor, hadronic, and cosmological observables.  The
construction uses the defined information constant
$\kappa=\ln 2/(24\pi)$, three discrete integers $(L_3,L_4,L_5)=(7,2,5)$,
prime-valued flavor labels, and a small set of external physical anchors.
The manuscript does not claim that these ingredients have been derived from
an accepted microscopic theory.  Instead, it distinguishes established
mathematics, model postulates, retrospective formula assignments, diagnostic
no-go results, and prospectively frozen tests.

The retrospective ledger contains close numerical matches in several sectors,
but these matches are not combined into a single global accuracy score because
the observables have different units, uncertainties, dependence structures,
and degrees of post-selection.  Two failures are kept explicit: the quoted
gauge-boson masses remain approximately $4.5$--$5.0\%$ high, and the fixed
strong-CP-to-neutron-EDM mapping gives
$d_n=5.33\times10^{-26}\,e\,\mathrm{cm}$, a factor $2.96$ above the current
$90\%$ confidence upper bound used in this work.  The Collatz $3/4$ mass trace
is also retained as an incomplete comparative signal rather than a closed mass
derivation.  A restricted common-baseline repair is excluded by an explicit
no-go theorem.  Three separable candidates are frozen for future direct
Higgs-coupling likelihoods, with no candidate deletion or formula alteration
permitted after inspection of the named data.

The complete source, validation ledgers, and reproducible LaTeX build are
available in the accompanying repository.

\section*{Keywords}
geometric phase; information geometry; phenomenological ansatz; flavor;
Standard Model observables; Collatz dynamics; Ramanujan structure;
preregistration; reproducibility.

\section*{Claim and evidence taxonomy}
\begin{description}
  \item[A -- established result:] a standard mathematical identity or external
  experimental result supported by the cited literature.
  \item[B -- model postulate:] a defining structural choice of Metatime/TIR.
  \item[C -- retrospective assignment:] a formula assessed after some or all of
  the comparison data were known.
  \item[D -- diagnostic or no-go:] an internal audit, failure, or restricted
  impossibility result.
  \item[E -- prospective prediction:] a frozen formula, observable, and decision
  rule specified before the named future dataset is inspected.
  \item[F -- external anchor:] a measured scale or coefficient supplied to the
  construction rather than predicted by it.
\end{description}

\section*{Availability and version status}
This is Version 11.0, designated a \emph{Publication Candidate}.  It is suitable
for circulation as a research monograph or preprint.  It is not presented as an
experimentally established derivation of the Standard Model from first
principles.  Source code, version history, machine-readable validation outputs,
and the build workflow are maintained at
\url{https://github.com/AdrianLipa90/The-Fundamental-Theory-of-Informational-Relations}.

\section*{Declarations}
This manuscript is an independent research work.  No external grant is
identified in the repository record.  No competing interest is declared in
this manuscript version.

\cleardoublepage
